% GENERATED FILE. Edit canonical lessons, then run npm run generate:books.
% canonical-source-hash: fnv1a64:023e1f16565a08a2
% canonical-lessons: TA-C34-edu, TA-C34-kel, TA-W14-read-pesu, TA-C34-utavu, TA-C34-pidi

\chapter{Four Verbs Between People}
\label{ch:ta-verbs-between-people}
\hlchaptermodality{34}

\begin{chapteropening}
\textbf{By the end of this chapter:} \emph{I can say I take, I ask, I help and I like in Tamil, hear \ta{கேள்} doing the work of both asking and hearing, choose between \ta{பிடிக்கும்} and \ta{விரும்பு}, and say \ta{எனக்குப்} \ta{பிடிக்கும்} knowing why the liker sits in the dative exactly as with \ta{தெரியும்}.}

Say \ta{எனக்குத்} \ta{தமிழ்} \ta{பிடிக்கும்} and unpack it as \textquotedblleft{}to me, Tamil catches\textquotedblright{}; name the three dative-experiencer verbs \ta{தெரியும்} / \ta{புரிகிறது} / \ta{பிடிக்கும்}, contrast \ta{விரும்பு}’s ordinary subject, and place Kannada \ka{ಹಿಡಿ} on the p$\to$h shift met in Chapter 7.
\end{chapteropening}

\section[eḍu]{\ta{எடு} (eḍu) — \textquotedblleft{}to take,\textquotedblright{} beside the \textquotedblleft{}put\textquotedblright{} you already own}

\label{lesson:TA-C34-edu}



\begin{quote}
The last chapter's verbs happened inside your head. These four happen between you and somebody else. Say \emph{paḍikkiṟēṉ} and \emph{eḻudugiṟēṉ} once more, then put your hands back in the world.
\end{quote}

\subsection*{You'll want to know: \ta{எடு}}
\begin{quote}
\textbf{\ta{எடு}} — \emph{eḍu} — \textbf{to take; to pick up, to lift, to remove}
\end{quote}

\textbf{\ta{எடுக்கிறேன்}} (\emph{eḍukkiṟēṉ}), \textquotedblleft{}I take.\textquotedblright{} The \emph{kk} is held, so \textbf{\ta{எடு} is strong}; past \textbf{\ta{எடுத்தேன்}} (\emph{eḍuttēṉ}). Two sentences you can already build in full:

\begin{quote}
\textbf{\ta{நான்} \ta{தண்ணீர்} \ta{எடுக்கிறேன்}.} — \emph{nāṉ taṇṇīr eḍukkiṟēṉ} — \textquotedblleft{}I take water.\textquotedblright{}
\end{quote}

>

\begin{quote}
\textbf{\ta{நான்} \ta{அரிசி} \ta{எடுக்கிறேன்}.} — \emph{nāṉ arisi eḍukkiṟēṉ} — \textquotedblleft{}I take rice.\textquotedblright{}
\end{quote}

\subsection*{The letters in this word}
\textbf{\ta{எ}} is the standing vowel \emph{e}, from \textbf{\ta{எழுது}}; \textbf{\ta{டு}} is \textbf{\ta{ட}} with the \emph{u} sign, and \textbf{\ta{ட}} came apart in \textbf{\ta{படி}}. Nothing here is new — the first time in this book that has been true of a whole headword. The \emph{ḍ} you hear is the same positional softening: \textbf{\ta{ட}} goes soft between vowels and hard when doubled, as in \emph{eḍuttēṉ}.

\begin{cousinweb}
You have met this verb's near-twin already, hiding inside a meal. \textbf{\ta{சாப்பிடு}} (\emph{sāppiḍu}), \textquotedblleft{}to eat,\textquotedblright{} is \emph{sāppu} (food) plus \textbf{\ta{இடு}} (\emph{iḍu}), \textquotedblleft{}\textbf{to put}\textquotedblright{} — the light verb that lends its endings to a noun.

\noindent
\begin{tabularx}{\linewidth}{@{}>{\raggedright\arraybackslash}X>{\raggedright\arraybackslash}X@{}}
\toprule
\textbf{verb} & \textbf{what it does} \\
\midrule
\textbf{\ta{இடு}} \emph{iḍu} & put it \textbf{down} \\
\textbf{\ta{எடு}} \emph{eḍu} & pick it \textbf{up} \\
\bottomrule
\end{tabularx}

One standing vowel apart — \textbf{\ta{இ}} against \textbf{\ta{எ}} — pointing opposite ways. That is a gift to your memory and \textbf{not} an etymology: Tamil's dictionaries keep them as \textbf{separate roots}, and a tidy pair of opposites is exactly the resemblance that fools people. Use it to remember; never as evidence.

What \emph{is} inherited is \emph{eḍu} itself. Malayalam has \textbf{\ml{എടുക്കുക}} (\emph{eṭukkuka}), the same verb; Kannada and Telugu say \textbf{\ka{ಎತ್ತು}} / \textbf{\te{ఎత్తు}} (\emph{ettu}), \textquotedblleft{}to lift,\textquotedblright{} from the same old root with a doubled \emph{t} where Tamil keeps its retroflex.
\end{cousinweb}

\subsection*{Your turn}
Say these aloud:

\begin{itemize}
\raggedright
  \item the opposed pair — \textquotedblleft{}iḍu, put … eḍu, take\textquotedblright{}
  \item \textquotedblleft{}nāṉ taṇṇīr eḍukkiṟēṉ\textquotedblright{} — I take water
  \item \textquotedblleft{}nāṉ arisi eḍukkiṟēṉ\textquotedblright{} — I take rice
  \item the cousins lifting — \textquotedblleft{}eṭukkuka … ettu … ettu\textquotedblright{}
\end{itemize}

\begin{tcolorbox}[breakable,colback=teal!4,colframe=teal!35!black,title={Before you move on}]
What does \emph{eḍu} mean, and which camp is it in? (\textbf{To take}; \textbf{strong} — \emph{eḍukkiṟēṉ} holds the \emph{kk}.) Which verb is one vowel away, and where did you meet it? (\emph{\textbf{\ta{இடு}}}, \textquotedblleft{}to put\textquotedblright{} — inside \textbf{\ta{சாப்பிடு}}.) Is that pair related? (\textbf{No} — the dictionaries keep them apart; it is a memory hook, not evidence.) Say \textquotedblleft{}I take water,\textquotedblright{} then \textquotedblleft{}I take rice.\textquotedblright{} (\emph{Nāṉ taṇṇīr eḍukkiṟēṉ}; \emph{nāṉ arisi eḍukkiṟēṉ}.)
\end{tcolorbox}
\section[kēḷ]{\ta{கேள்} (kēḷ) — one verb for asking and for hearing}

\label{lesson:TA-C34-kel}



\begin{quote}
\emph{Eḍukkiṟēṉ} — I take. Name the verb it is one vowel away from. Now a verb doing two English jobs at once.
\end{quote}

\subsection*{You'll want to know: \ta{கேள்}}
\begin{quote}
\textbf{\ta{கேள்}} — \emph{kēḷ} — \textbf{to ask; to hear, to listen}
\end{quote}

One syllable, ending on retroflex \textbf{\ta{ள்}}. The everyday form surprises:

\begin{quote}
\textbf{\ta{நான்} \ta{கேட்கிறேன்}.} — \emph{nāṉ kēṭkiṟēṉ} — \textquotedblleft{}I ask.\textquotedblright{} (Also: \textquotedblleft{}I hear.\textquotedblright{})
\end{quote}

The \textbf{\ta{ள்}} has hardened to \textbf{\ta{ட்}} — \textbf{\ta{வா}}'s trick a second time: you say it \emph{vā}, but the endings attach to \emph{varu-}. Here they attach to \textbf{\ta{கேட்}-} (\emph{kēṭ-}), and the past is \textbf{\ta{கேட்டேன்}} (\emph{kēṭṭēṉ}), a strong verb's doubled consonant. Two verbs with two stems apiece.

\subsection*{The letters in this word}
\textbf{\ta{கே}} is \textbf{\ta{க}} with the long-\emph{ē} sign \textbf{\ta{ே}}; \textbf{\ta{ள்}} is retroflex \textbf{\ta{ள}} shut off by a \textbf{puḷḷi} — the dot Chapter 8 pointed at inside \textbf{\ta{செய்து}}. In \textbf{\ta{கேட்கிறேன்}} the seam is \textbf{\ta{ட்க}}.

\begin{cousinweb}
English keeps \emph{ask} and \emph{hear} apart. Tamil does not, and neither does its noun: \textbf{\ta{கேள்வி}} (\emph{kēḷvi}) is \textbf{\textquotedblleft{}a question\textquotedblright{}} and also \textquotedblleft{}a hearing.\textquotedblright{} The family splits down the middle:

\begin{itemize}
\raggedright
  \item Kannada \textbf{\ka{ಕೇಳು}} (\emph{kēḷu}) does \textbf{both} jobs, as Tamil does.
  \item Malayalam \textbf{\ml{കേൾക്കുക}} (\emph{kēḷkkuka}) kept only \textquotedblleft{}hear.\textquotedblright{}
  \item Telugu splits them: \textbf{\te{విను}} (\emph{vinu}), hear; \textbf{\te{అడుగు}} (\emph{aḍugu}), ask.
\end{itemize}

The verb is inherited, but four languages cannot settle which arrangement came first; treat the fusion as Tamil's and Kannada's shared habit alone.
\end{cousinweb}

\begin{culture}
Tamil is \textbf{diglossic}: its literary variety and the variety people speak have drifted far apart. This verb makes that impossible to ignore.

\begin{itemize}
\raggedright
  \item \textbf{\ta{கேள்}} is the dictionary and literary form.
  \item \textbf{\ta{கேட்கிறேன்}} is standard spoken Tamil — \textbf{the register these lessons teach}, as every lesson's variety line has said.
  \item Relaxed speech says \textbf{\ta{கேக்கிறேன்}} (\emph{kēkkiṟēṉ}) and \textbf{\ta{கேக்க}} (\emph{kēkka}). Recognise them; you need not say them.
\end{itemize}

For a polite request use Chapter 8's respectful ending: \textbf{\ta{கேளுங்கள்}} (\emph{kēḷungaḷ}), \textquotedblleft{}please ask.\textquotedblright{} Keep \textbf{\ta{தயவுசெய்து}} for emphasis.
\end{culture}

\subsection*{Your turn}
Say these aloud:

\begin{itemize}
\raggedright
  \item the two shifting stems — \textquotedblleft{}vā … varugiṟēṉ\textquotedblright{} then \textquotedblleft{}kēḷ … kēṭkiṟēṉ\textquotedblright{}
  \item \textquotedblleft{}kēḷvi\textquotedblright{} — a question — then the polite \textquotedblleft{}kēḷungaḷ\textquotedblright{}
  \item the family split — \textquotedblleft{}Kannada kēḷu, both … Telugu vinu and aḍugu, two\textquotedblright{}
\end{itemize}

\begin{tcolorbox}[breakable,colback=teal!4,colframe=teal!35!black,title={Before you move on}]
What two English verbs is \emph{kēḷ}, and what happens to its \textbf{\ta{ள்}} before an ending? (\textbf{Ask} and \textbf{hear}; it hardens to \textbf{\ta{ட்}}.) Which verb taught you to expect a second stem, and what is \emph{kēḷvi}? (\emph{\textbf{\ta{வா}}}, built on \emph{varu-}; \textbf{a question}.) Which register do these lessons teach, and how do you ask politely? (\textbf{Standard spoken} \emph{kēṭkiṟēṉ}; \emph{\textbf{kēḷungaḷ}}.)
\end{tcolorbox}
\section[pēsu]{\ta{பேசு} — the long partner of \ta{ெ}}

\label{lesson:TA-W14-read-pesu}



\begin{quote}
Twice now a Tamil vowel has come in a short and a long version: \textbf{\ta{அ}} and \textbf{\ta{ஆ}}, \textbf{\ta{ி}} and \textbf{\ta{ீ}}. For \textbf{\ta{அ}}/\textbf{\ta{ஆ}} the script data even says the long one is the short one plus a tail; for \textbf{\ta{ி}}/\textbf{\ta{ீ}} it says nothing, and the pairing is the lesson's. Here is a third pair.
\end{quote}

\subsection*{Script you'll notice: the \ta{ே} sign}
\textbf{\ta{ே}} is the sign for long \emph{ē}, and it is the long partner of the \textbf{\ta{ெ}} you met in \textbf{\ta{பெயர்}}. Like \textbf{\ta{ெ}}, it is written to the \textbf{left} of its consonant and spoken after it.

\begin{quote}
\textbf{\ta{ப} + \ta{ே} $\to$ \ta{பே}}, said \emph{pē}.
\end{quote}

\noindent
\begin{tabularx}{\linewidth}{@{}>{\raggedright\arraybackslash}X>{\raggedright\arraybackslash}X>{\raggedright\arraybackslash}X@{}}
\toprule
\textbf{sign} & \textbf{sound} & \textbf{where it is written} \\
\midrule
\textbf{\ta{ெ}} & short \emph{e} & to the left \\
\textbf{\ta{ே}} & long \emph{ē} & to the left \\
\textbf{\ta{ை}} & \emph{ai} & to the left \\
\bottomrule
\end{tabularx}

Three left-standing signs now. The evidence is the same as it was for \textbf{\ta{ெ}}: the script data describes \textbf{\ta{ை}} and has no entry for \textbf{\ta{ெ}} or \textbf{\ta{ே}} at all, so both the shared position \textbf{and} the short/long pairing are stated from the pattern these signs share, not from a source for either one.

\subsection*{Script you'll notice: build the word}
\noindent
\begin{tabularx}{\linewidth}{@{}>{\raggedright\arraybackslash}X>{\raggedright\arraybackslash}X>{\raggedright\arraybackslash}X@{}}
\toprule
\textbf{piece} & \textbf{} & \textbf{} \\
\midrule
\textbf{\ta{பே}} & \emph{pē} — \ta{ப} with the \ta{ே} sign on its left & \ta{ப} from \textbf{\ta{பெயர்}} \\
\textbf{\ta{சு}} & \emph{su} — \ta{ச} with the \ta{ு} sign & \ta{ச} from \textbf{\ta{சரி}}, \ta{ு} from \textbf{\ta{இருக்கிறீர்கள்}} \\
\bottomrule
\end{tabularx}

\begin{quote}
\textbf{\ta{பே} · \ta{சு} $\to$ \ta{பேசு}}
\end{quote}

Two pieces, one new sign, and its two signs pull in opposite directions on the page: \textbf{\ta{ே}} stands before its letter, \textbf{\ta{ு}} joins after it. \textbf{\ta{பேசு}} is \textquotedblleft{}speak\textquotedblright{} — the verb every lesson so far has been in service of.

\subsection*{Your turn}
\begin{itemize}
\raggedright
  \item \emph{Read it:} \textbf{\ta{பே}} — then \textbf{\ta{பேசு}}
  \item \emph{Say it:} \textquotedblleft{}pēsu\textquotedblright{} — \textquotedblleft{}speak\textquotedblright{}
  \item \emph{Contrast:} \textbf{\ta{பெ}} and \textbf{\ta{பே}} — the short sign and the long one, both on the left
  \item \emph{Contrast:} \textbf{\ta{பே}} and \textbf{\ta{சு}} — a sign before, a sign after, inside one word
  \item \emph{Read it:} \textbf{\ta{தமிழ்} \ta{பேசு}} — \textquotedblleft{}speak Tamil,\textquotedblright{} both words off the page
\end{itemize}

\begin{tcolorbox}[breakable,colback=teal!4,colframe=teal!35!black,title={Before you move on}]
Read \textbf{\ta{பேசு}}. Which side of its consonant is \textbf{\ta{ே}} written on? (\textbf{The left} — like \textbf{\ta{ெ}} and \textbf{\ta{ை}}.) Which sign is its short partner? (\textbf{\ta{ெ}}, as in \textbf{\ta{பெயர்}}.) What is the \textbf{\ta{ச}} doing here — is it \emph{s} or \emph{ch}? (\textbf{Either}; one letter covers both, and position picks.) Next: the verb for helping.
\end{tcolorbox}
\section[udavu]{\ta{உதவு} (udavu) — \textquotedblleft{}to help,\textquotedblright{} and a word Tamil did not borrow}

\label{lesson:TA-C34-utavu}



\begin{quote}
\emph{Kēṭkiṟēṉ} — I ask. \emph{Kēḷungaḷ} — please ask. Say both, and keep that polite ending in your ear; it is about to arrive on a second verb unchanged.
\end{quote}

\subsection*{You'll want to know: \ta{உதவு}}
\begin{quote}
\textbf{\ta{உதவு}} — \emph{udavu} — \textbf{to help; to be of use}
\end{quote}

\textbf{\ta{உதவுகிறேன்}} (\emph{udavugiṟēṉ}), \textquotedblleft{}I help.\textquotedblright{} The single softened \emph{-giṟ-} puts \textbf{\ta{உதவு}} in the \textbf{weak} camp with \emph{pō} and \emph{eḻudu}, and the noun is one step away: \textbf{\ta{உதவி}} (\emph{udavi}), \textquotedblleft{}\textbf{help}.\textquotedblright{}

\begin{quote}
\textbf{\ta{நான்} \ta{உங்களுக்கு} \ta{உதவுகிறேன்}.} — \emph{nāṉ ungaḷukku udavugiṟēṉ} — \textquotedblleft{}I help you.\textquotedblright{}
\end{quote}

The person helped takes the dative \textbf{-\ta{உக்கு}} from Chapter 6, and the polite form needs nothing new: \textbf{\ta{உதவுங்கள்}} (\emph{udavungaḷ}), \textquotedblleft{}please help\textquotedblright{} — the same respectful \textbf{-\ta{உங்கள்}} you just put on \emph{kēḷ}.

\subsection*{The letters in this word}
\textbf{\ta{உ}} is the standing vowel \emph{u}, used when a word begins with it; then \textbf{\ta{த}} and \textbf{\ta{வு}}, which is \textbf{\ta{வ}} with the \emph{u} sign. The \emph{d} you hear is \textbf{\ta{த}} between two vowels, softening exactly as \textbf{\ta{ட}} did in \emph{paḍi} and \textbf{\ta{க}} does in \emph{pōgiṟēṉ}.

\begin{cousinweb}
Tamil turns verbs into nouns with a small ending, and you have collected three without being told to: \textbf{\ta{புரி}} $\to$ \textbf{\ta{புரிதல்}} (\emph{puridal}), understanding; \textbf{\ta{கேள்}} $\to$ \textbf{\ta{கேள்வி}} (\emph{kēḷvi}), a question; \textbf{\ta{உதவு}} $\to$ \textbf{\ta{உதவி}} (\emph{udavi}), help.

Now the part worth keeping. For \textquotedblleft{}help,\textquotedblright{} the neighbours nearly all reach for Sanskrit \textbf{\dv{सहाय}} (\emph{sahāya}): Telugu \textbf{\te{సహాయం}}, Kannada \textbf{\ka{ಸಹಾಯ}}, Malayalam \textbf{\ml{സഹായം}}, Hindi \textbf{\dv{सहायता}}. Tamil's everyday word is \textbf{\ta{உதவி}}, its own.

Hold that against Chapter 8, where \emph{please} went the other way: \textbf{all four} sisters built it on Sanskrit \textbf{\dv{दया}} (\emph{daya}) plus a light verb — \emph{tayavuseytu}, \emph{dayaviṭṭu}, \emph{dayacēsi}, \emph{dayavāyi}. Same four languages, opposite outcome. So the rule is not \textquotedblleft{}Tamil refuses to borrow\textquotedblright{}; it is that \textbf{borrowing happens word by word}, and the only way to know is to check each one.
\end{cousinweb}

\subsection*{Your turn}
Say these aloud:

\begin{itemize}
\raggedright
  \item \textquotedblleft{}nāṉ ungaḷukku udavugiṟēṉ\textquotedblright{} — I help you
  \item the two polite forms — \textquotedblleft{}kēḷungaḷ … udavungaḷ\textquotedblright{}
  \item the three nouns — \textquotedblleft{}puridal … kēḷvi … udavi\textquotedblright{}
  \item the borrowed one and the kept one — \textquotedblleft{}tayavuseytu … udavi\textquotedblright{}
\end{itemize}

\begin{tcolorbox}[breakable,colback=teal!4,colframe=teal!35!black,title={Before you move on}]
Which camp is \emph{udavu} in, and what is \emph{udavi}? (\textbf{Weak}, by the single softened bead; \textbf{help}, the noun.) Say \textquotedblleft{}I help you,\textquotedblright{} then \textquotedblleft{}please help.\textquotedblright{} (\emph{Nāṉ ungaḷukku udavugiṟēṉ}; \emph{udavungaḷ}.) Which word do the neighbours borrow from Sanskrit where Tamil keeps its own? (\textbf{Help} — \emph{sahāya} against \emph{udavi}.) And which word did all four sisters borrow together? (\emph{\textbf{Please}} — \emph{daya} plus a light verb.)
\end{tcolorbox}
\section[piḍi]{\ta{பிடி} (piḍi) — \textquotedblleft{}to catch,\textquotedblright{} and therefore \textquotedblleft{}to like\textquotedblright{}}

\label{lesson:TA-C34-pidi}



\begin{quote}
Say this chapter's three verbs — \emph{eḍukkiṟēṉ}, \emph{kēṭkiṟēṉ}, \emph{udavugiṟēṉ} — then the two sentences with no \emph{nāṉ}: \emph{enakku teriyum}, \emph{enakku purigiṟadu}.
\end{quote}

\subsection*{You'll want to know: \ta{பிடி}}
\begin{quote}
\textbf{\ta{பிடி}} — \emph{piḍi} — \textbf{to catch, to seize, to hold}
\end{quote}

\textbf{\ta{பிடிக்கிறேன்}} (\emph{piḍikkiṟēṉ}) is \textquotedblleft{}I catch, I hold,\textquotedblright{} and the held \emph{kk} puts it in the \textbf{strong} camp with \emph{ninai}, \emph{paḍi} and \emph{eḍu}.

\begin{grammarlens}[title={Grammar Lens: \ta{எனக்குப்} \ta{பிடிக்கும்}}]
\begin{quote}
\textbf{\ta{எனக்குத்} \ta{தமிழ்} \ta{பிடிக்கும்}.} — \emph{enakkut tamiḻ piḍikkum} — \textquotedblleft{}I like Tamil.\textquotedblright{}
\end{quote}

Word for word: \textquotedblleft{}\textbf{to me}, Tamil catches.\textquotedblright{} The Tamil does the catching; you are caught. \textbf{\ta{பிடிக்கும்}} carries \textbf{no person at all} — the person-less ending of \textbf{\ta{தெரியும்}} — so only the dative word changes: \emph{enakku piḍikkum}, \emph{uṉakku piḍikkum}, \emph{avaṉukku piḍikkum}.

That makes three: \textbf{\ta{தெரியும்}}, \textbf{\ta{புரிகிறது}}, \textbf{\ta{பிடிக்கும்}} — knowing, understanding, liking. Three endings; each puts the experiencer in the dative and the thing in the subject. Chapter 6's \textbf{dative subject} was never an exception.
\end{grammarlens}

\subsection*{You'll want to know: \ta{விரும்பு}}
\begin{quote}
\textbf{\ta{விரும்பு}} — \emph{virumbu} — \textbf{to want, to like, to love}
\end{quote}

An ordinary verb with an ordinary subject: \textbf{\ta{நான்} \ta{விரும்புகிறேன்}} (\emph{nāṉ virumbugiṟēṉ}), \textquotedblleft{}I want.\textquotedblright{} Weak; its noun is \textbf{\ta{விருப்பம்}} (\emph{viruppam}), a wish. Not a synonym: \textbf{\ta{பிடிக்கும்}} is taste, which you had no say in; \textbf{\ta{விரும்பு}} is choice.

\begin{cousinweb}
The grasp-verb is old and shared: Malayalam \textbf{\ml{പിടിക്കുക}} (\emph{piṭikkuka}) and Telugu \textbf{\te{పట్టు}} (\emph{paṭṭu}) both mean \textquotedblleft{}catch, hold.\textquotedblright{} Kannada's is \textbf{\ka{ಹಿಡಿ}} (\emph{hiḍi}) — not a different word: Old Kannada wrote \emph{piḍi}, and the \textbf{p $\to$ h} shift behind \emph{pattu / hattu} and \emph{pāl / hālu} turned it into \emph{hiḍi}. A third instance of a rule you have already met.

The wider point: Spanish says \emph{me gusta}, Italian \emph{mi piace}, Hindi \emph{mujhe pasand hai} — each putting the liker in the dative and the liked thing in the subject, exactly as \emph{enakkut tamiḻ piḍikkum} does, and none of them related to Tamil. \textbf{Two unrelated families reached the same design independently.}
\end{cousinweb}

\subsection*{Your turn}
Say these aloud:

\begin{itemize}
\raggedright
  \item the three dative sentences — \textquotedblleft{}enakku teriyum … purigiṟadu … piḍikkum\textquotedblright{}
  \item the other design — \textquotedblleft{}nāṉ virumbugiṟēṉ\textquotedblright{}, I want
  \item the chapter — \textquotedblleft{}eḍukkiṟēṉ, kēṭkiṟēṉ, udavugiṟēṉ, piḍikkum\textquotedblright{}
  \item the Kannada shift — \textquotedblleft{}piḍi … hiḍi\textquotedblright{}, as \textquotedblleft{}pattu … hattu\textquotedblright{}
\end{itemize}

\begin{tcolorbox}[breakable,colback=teal!4,colframe=teal!35!black,title={Before you move on}]
Say \textquotedblleft{}I like Tamil\textquotedblright{} and name its subject. (\emph{Enakkut tamiḻ piḍikkum}, \textquotedblleft{}to me, Tamil catches\textquotedblright{}; the subject is \textbf{Tamil}.) Name the three verbs built that way, and what \emph{virumbu} does differently. (\emph{\textbf{Teriyum}}, \emph{\textbf{purigiṟadu}}, \emph{\textbf{piḍikkum}}; \emph{virumbu} takes an \textbf{ordinary subject}.) Why is Kannada's \emph{hiḍi} the same word as \emph{piḍi}, and which unrelated language builds \textquotedblleft{}I like\textquotedblright{} this way? (\textbf{The p $\to$ h shift}; \textbf{Spanish} \emph{me gusta}.)
\end{tcolorbox}
